%% ============================================================================
%% Section: Indexing and Query Processing (CORE)
%% Owns Table tab:complexity and Fig fig:index.
%% Design follows CODE-PLAN.md WP1 (layout, closures) and WP2 (algorithms).
%% ============================================================================
\section{Indexing and Query Processing}\label{sec:index}

This section is the core of the paper. We present the physical layout the
snapshot compiles to (Section~\ref{sec:index:layout}), the precomputed
reachability closures that accelerate every structural query
(Section~\ref{sec:index:bitsets}), and the query-processing algorithms built
on them: reachability and path roles (Section~\ref{sec:index:reach}),
conditioning-aware d-separation (Section~\ref{sec:index:dsep}), and
identification by adjustment (Section~\ref{sec:index:adj}).
Table~\ref{tab:complexity} summarizes the per-operator cost against a
map-of-sets baseline, distinguishing the operators that become sublinear from
those that keep a linear traversal with a smaller constant.

\subsection{Physical layout}\label{sec:index:layout}

The snapshot's variables are assigned dense integer indices $0..n-1$ through a
\textsf{name}$\to$\textsf{index} map, so that every downstream structure is an
array rather than a hash map (Fig.~\ref{fig:index}). Edges are stored as
compressed sparse rows: a forward and a reverse offset/target array pair, so
that the children of $i$ are the slice
$\textsf{fwdTarget}[\textsf{fwdOffset}[i] .. \textsf{fwdOffset}[i{+}1])$ and
the parents of $i$ the symmetric slice of the reverse pair, giving
cache-friendly, allocation-free neighbor iteration. The topological order is
stored as an \texttt{int[]} together with its inverse, and each node carries
its layer $\textsf{layer}[i]$ (0 for sources, otherwise one more than the
maximum layer of its parents), which lets a traversal prune any neighbor that
cannot lie on a path between two nodes by a layer comparison. This replaces
the baseline representation, a \lstinline[style=causaldsl]|Map<String,Set<String>>|
adjacency with string-keyed lookups and per-query depth-first search, whose
every structural query costs an $O(|V|+|E|)$ traversal with hashing and
boxing overhead.

\subsection{Reachability closures as bitsets}\label{sec:index:bitsets}

On top of the dense layout we precompute, for every variable, its ancestor set
and its descendant set as fixed-width bitsets
$\textsf{anc}[i], \textsf{des}[i] \in \{0,1\}^n$ packed into
$w = \lceil n/64 \rceil$ machine words. The ancestor closures are filled in a
single forward-topological sweep with the recurrence
\[
  \textsf{anc}[i] \;=\; \bigcup_{p \in \pa(i)} \bigl(\textsf{anc}[p] \cup \{p\}\bigr),
\]
and the descendant closures symmetrically in reverse-topological order. Each
sweep performs one word-parallel union per edge, so the build costs
$O((|V|+|E|)\, w)$ time and the two closures occupy $2\,|V|\,w$ words of
space, i.e., $O(|V|^2/64)$. This is a deliberate space-for-time trade: it
turns reachability from a traversal into a bit test and, more importantly for
the causal workload, delivers whole ancestor and descendant sets as single
array reads. We discuss when the trade stops paying at the end of this
section.

\begin{figure}[t]\centering
  \resizebox{\linewidth}{!}{%% ============================================================================
%% Figure: compiled index layout (dense CSR + reachability bitsets)
%% Included from sections/indexing.tex inside a figure env via \resizebox{\input}.
%% Requires: \usepackage{tikz}; \usetikzlibrary{positioning,fit}
%% Example DAG: A(0)->B(1), A(0)->C(2), B(1)->D(3), C(2)->D(3)
%% ============================================================================
\begin{tikzpicture}[
    font=\scriptsize,
    >={Stealth[length=1.6mm]},
    cell/.style={draw, minimum width=5.2mm, minimum height=5mm, inner sep=1pt,
                 align=center, fill=white},
    lab/.style={anchor=east, font=\scriptsize\itshape},
    node distance=2mm,
]

% ---- tiny example DAG (top-left) ----
\node[draw, circle, minimum size=5mm] (A) at (0,0)      {A};
\node[draw, circle, minimum size=5mm] (B) at (1.1,0.7)  {B};
\node[draw, circle, minimum size=5mm] (C) at (1.1,-0.7) {C};
\node[draw, circle, minimum size=5mm] (D) at (2.2,0)    {D};
\draw[->] (A) -- (B); \draw[->] (A) -- (C);
\draw[->] (B) -- (D); \draw[->] (C) -- (D);
\node[font=\scriptsize\itshape] at (1.1,-1.6) {DAG (indices A{=}0,B{=}1,C{=}2,D{=}3)};

% ---- name -> index map ----
\node[lab] (m0) at (5.2, 1.0) {nameToIndex:};
\node[cell] (mA) at (6.1,1.0) {A:0};
\node[cell, right=of mA] (mB) {B:1};
\node[cell, right=of mB] (mC) {C:2};
\node[cell, right=of mC] (mD) {D:3};

% ---- forward CSR ----
\node[lab] (f0) at (5.2, 0.3) {forwardAdj:};
\node[cell] (f1) at (6.1,0.3) {1,2};
\node[cell, right=of f1] (f2) {3};
\node[cell, right=of f2] (f3) {3};
\node[cell, right=of f3] (f4) {--};

% ---- reverse CSR ----
\node[lab] (r0) at (5.2, -0.4) {reverseAdj:};
\node[cell] (r1) at (6.1,-0.4) {--};
\node[cell, right=of r1] (r2) {0};
\node[cell, right=of r2] (r3) {0};
\node[cell, right=of r3] (r4) {1,2};

% ---- topo order / layer ----
\node[lab] (t0) at (5.2, -1.1) {topoOrder:};
\node[cell] (t1) at (6.1,-1.1) {0};
\node[cell, right=of t1] (t2) {1};
\node[cell, right=of t2] (t3) {2};
\node[cell, right=of t3] (t4) {3};
\node[lab] (l0) at (5.2, -1.7) {topoLayer:};
\node[cell] (l1) at (6.1,-1.7) {0};
\node[cell, right=of l1] (l2) {1};
\node[cell, right=of l2] (l3) {1};
\node[cell, right=of l3] (l4) {2};

% ---- ancestor bitsets (rows = variable, columns = bit for A B C D) ----
\node[font=\scriptsize\itshape] at (1.4,-2.6) {ancestorBits[i] (bit set = ancestor)};
\node[lab] at (-0.3,-3.2) {A}; \node[lab] at (-0.3,-3.7) {B};
\node[lab] at (-0.3,-4.2) {C}; \node[lab] at (-0.3,-4.7) {D};
% rows explicitly (A,B,C,D) x (A,B,C,D)
\foreach \row/\y/\vals in {A/-3.2/{0/0,1/0,2/0,3/0}, B/-3.7/{0/1,1/0,2/0,3/0},
                           C/-4.2/{0/1,1/0,2/0,3/0}, D/-4.7/{0/1,1/1,2/1,3/0}} {
  \foreach \c/\v in \vals {
    \node[cell, minimum width=4.5mm] at ({0.1+\c*0.5},\y) {\v};
  }
}

% ---- descendant bitsets ----
\node[font=\scriptsize\itshape] at (5.2,-2.6) {descendantBits[i] (bit set = descendant)};
\node[lab] at (3.6,-3.2) {A}; \node[lab] at (3.6,-3.7) {B};
\node[lab] at (3.6,-4.2) {C}; \node[lab] at (3.6,-4.7) {D};
\foreach \row/\y/\vals in {A/-3.2/{0/0,1/1,2/1,3/1}, B/-3.7/{0/0,1/0,2/0,3/1},
                           C/-4.2/{0/0,1/0,2/0,3/1}, D/-4.7/{0/0,1/0,2/0,3/0}} {
  \foreach \c/\v in \vals {
    \node[cell, minimum width=4.5mm] at ({4.0+\c*0.5},\y) {\v};
  }
}

\end{tikzpicture}
}
  \caption{Compiled index layout for a small fragment. Variables map to dense
  indices; forward and reverse CSR arrays store adjacency; a topological order
  and a per-node layer drive pruning; per-node ancestor and descendant bitsets
  ($\lceil n/64\rceil$ words each) are filled in forward- and
  reverse-topological order.}
  \label{fig:index}
\end{figure}

\subsection{Reachability and path-role queries}\label{sec:index:reach}

With the closures in hand the reachability operators are immediate: $\an(X)$
and $\de(X)$ are one array read, and a directed path from $X$ to $Y$ exists if
and only if bit $Y$ is set in $\textsf{des}[X]$, an $O(1)$ test. The path-role
operators reduce to bitset algebra over the closures:
{\small
\begin{align*}
  \text{confounders}(T,Y) &= \textsf{anc}[T] \cap \textsf{anc}[Y], \\
  \text{mediators}(T,Y)   &= \textsf{des}[T] \cap \textsf{anc}[Y], \\
  \cn(T,Y)   &= \text{mediators}(T,Y) \cup \{Y\}, \\
  \forb(T,Y) &= \{T\} \cup \cn(T,Y) \cup \textstyle\bigcup_{C \in \cn(T,Y)} \textsf{des}[C], \\
  \adj(T,Y)  &= (\textsf{anc}[T] \cup \textsf{anc}[Y]) \\
             &\qquad \setminus (\forb(T,Y) \cup \{T,Y\}).
\end{align*}}
Each of the first three is $O(w)$ word operations; the forbidden set costs
one union per causal node, $O(|\cn(T,Y)|\, w)$, and the canonical set two
more word operations. Colliders between $X$ and $Y$ are the variables outside
$\textsf{anc}[X] \cup \textsf{anc}[Y]$ having one parent in
$\textsf{des}[X] \cup \{X\}$ and one in $\textsf{des}[Y] \cup \{Y\}$, which
is one bit test per parent of each candidate. The variables on backdoor paths
are found by a breadth-first search on the skeleton starting from $\pa(T)$
that never enters $\forb(T,Y)$, with the visited set as a bitset. In the
running example, $\text{confounders}(\textit{interventions},
\textit{reported\_cases})$ is a single intersection returning
\textit{covid19\_burden}, and $\forb$ contains \textit{mobility},
\textit{exposure\_sarscov2}, \textit{onset\_covid19}, and their descendants.

\subsection{Conditioning-aware d-separation}\label{sec:index:dsep}

D-separation is not reachability: adding a variable to the conditioning set
$Z$ can block one path and, when that variable is a collider or one of its
descendants, open another. The test must follow edge orientation and the role
of each node under $Z$, which we do with the Bayes-Ball
procedure~\cite{shachter1998bayesball}. A search propagates balls from $X$,
tracking for each visited node whether it was entered from a parent or from a
child; at each node the admissible continuations depend on whether the node is
in $Z$: a node in $Z$ entered from a parent bounces the ball back to its
parents and blocks it otherwise, while a node outside $Z$ passes a ball
entered from a parent to its children and a ball entered from a child to both
its parents and its children. $X$ and $Y$ are d-separated given $Z$ exactly
when no ball reaches $Y$. The procedure is conditioning-aware by construction
and visits each (node, direction) state at most once, so its cost is
$O(|V|+|E|)$ in the worst case. Bitsets do not change this bound. However,
they change the constant in two ways. First, $X$, $Y$, $Z$, the two visited
sets, and the set of nodes whose outgoing edges are cut are all bitsets, so
membership is one word test and the traversal allocates nothing. Second,
d-separation of $X$ and $Y$ given $Z$ depends only on the subgraph induced by
$\an(X \cup Y \cup Z) \cup X \cup Y \cup Z$~\cite{geiger1990identifying}, and
that set is one union of closures: the traversal is masked to it and never
enqueues a node outside, which on sparse DAGs with small $Z$ prunes most of
the graph. Section~\ref{sec:eval} measures the effect of masking separately
from the effect of the representation. A flag cutting the outgoing edges of
the treatment towards its causal nodes turns the same routine into the
separation test in the proper backdoor graph used next.

\subsection{Identification by adjustment}\label{sec:index:adj}

Identification asks whether the effect of $T$ on $Y$ is estimable by covariate
adjustment and, if so, for one or all minimal adjustment sets. A naive
approach enumerates candidate sets and validates each with a d-separation
test; since the number of candidates is exponential in $|V|$, no speed-up of
the inner test makes it scale. We instead build on the constructive results
of van der Zander et al.~\cite{vanderzander2014constructing,vanderzander2019separators},
whose steps consume exactly the sets the index stores.

\noindent\textbf{Decision.} The effect is identifiable by adjustment if and
only if the canonical set $\adj(T,Y)$ satisfies the adjustment criterion.
Since $\adj(T,Y)$ is disjoint from $\forb(T,Y)$ by construction, the decision
is one closure construction ($O(|\cn(T,Y)|\,w)$ word operations) followed by
one masked Bayes-Ball run in the proper backdoor graph. This is the default of
\lstinline[style=causaldsl]|IDENTIFY EFFECT|: a query with no open backdoor
path returns after a single traversal, and so does a query on a graph with
thousands of variables.

\noindent\textbf{Enumeration.} Minimal adjustment sets are minimal separators
of $T$ and $Y$ in the proper backdoor graph restricted to the allowed
variables $R = V \setminus (\forb(T,Y) \cup \{T,Y\})$. We implement the
separator framework of~\cite{vanderzander2019separators} in three routines.
(1) $\textsc{FindSep}(I,R)$ returns $\bigl(\an(\{T,Y\} \cup I) \cap R\bigr)
\setminus \{T,Y\}$ if that set separates $T$ and $Y$, and reports that no
separator containing $I$ within $R$ exists otherwise; the candidate is one
union of closures and the check one traversal.
(2) $\textsc{FindMinSep}(I,R)$ shrinks the result of \textsc{FindSep} to an
inclusion-minimal separator by two reachability passes in the moralized
ancestral graph of $T$, $Y$, and the candidate, keeping only the candidate
variables reachable from $T$ and then from $Y$; the ancestral graph is the
mask of the previous section and moralization is local to it.
(3) $\textsc{ListMinSep}(I,R)$ enumerates all minimal separators by
branching on one allowed variable $v$ at a time, either forcing it into $I$
or removing it from $R$, and pruning a branch as soon as \textsc{FindMinSep}
reports that no separator exists in it. Every leaf emits one minimal set and
every internal node performs a bounded number of closure unions and
traversals, so the enumeration runs with polynomial
delay~\cite{vanderzander2019separators}: the time between two consecutive
answers is polynomial in $|V|$ regardless of how many answers there are.
\lstinline[style=causaldsl]|RETURN ADJUSTMENT SETS| streams the result of
\textsc{ListMinSep} with an optional limit, and \textsc{FindMinSep} alone
answers a request for a single minimal set.

\begin{table}[t]\centering
\caption{Per-operator cost of the map-of-sets baseline with per-query DFS and
of the bitset index, with $w = \lceil |V|/64 \rceil$, $c = |\cn(T,Y)|$, and
$V_m, E_m$ the nodes and edges of the ancestral mask. Tier (a) operators are sublinear;
tier (b) keeps a linear traversal with bitset membership and ancestral
masking; tier (c) replaces exponential enumeration with a bounded number of
tier (a) and (b) steps per answer.}
\label{tab:complexity}
\footnotesize
\setlength{\tabcolsep}{3pt}
\begin{tabular}{@{}llll@{}}
\toprule
Operator & Baseline & Index & Tier \\
\midrule
ancestors / descendants & $O(|V|{+}|E|)$ & $O(w)$ & (a) \\
directed path test      & $O(|V|{+}|E|)$ & $O(1)$ & (a) \\
confounders / mediators & $O(|V|{+}|E|)$ & $O(w)$ & (a) \\
forbidden / canonical set & $O(c(|V|{+}|E|))$ & $O(c\,w)$ & (a) \\
colliders               & $O(|V|{+}|E|)$ & $O(|E|)$ bit tests & (a) \\
d-separation (masked)   & $O(|V|{+}|E|)$ & $O(|V_m|{+}|E_m|)$ & (b) \\
identifiability decision & $O(2^{|V|})$ tests & $O(c\,w)$ + 1 test & (c) \\
minimal sets (each)     & $O(2^{|V|})$ tests & poly. delay & (c) \\
\midrule
index space             & $O(|V|{+}|E|)$ & $O(|V|{+}|E| + |V|\,w)$ & \\
index build             & $O(|V|{+}|E|)$ & $O((|V|{+}|E|)\,w)$ & \\
\bottomrule
\end{tabular}
\end{table}

\noindent\textbf{When closures stop paying.} The closures cost $2|V|w$ words:
about 25~MB for $|V| = 10^4$ and about 2.5~GB for $|V| = 10^5$. For DAGs
beyond a configured memory budget the snapshot keeps only the CSR layout and
the layer array, and answers reachability by a layer-pruned depth-first search
that never follows a neighbor whose layer lies outside the interval between
the two endpoints; path-role and identification operators then run over
sets computed on demand, at the baseline cost of Table~\ref{tab:complexity}
minus the hashing overhead. The choice is made once at compile time from
$|V|$ and the budget, and Section~\ref{sec:eval} reports the latency and
memory of both branches so that the crossover is a measured quantity rather
than a guess. Causal DAGs curated by domain experts are, in our experience,
well below the crossover: the largest reference network we use has
$|V| = 1041$ and its closures fit in about 280~KB.
